% Source observation: analyses/018-2026-09-08-results-materials/observations/O012-manuscript-results.md
% Argument revision after adversarial review, 2026-09-05. Hand-edited.
% Sources and claim boundaries: positioning.md and reviews/2026-09-05-adversarial/.
% Result preview authorized 8 September 2026; evidence: Analysis 018 O002/O003/O007/O011.
\section{Introduction}
\label{sec:introduction}

Activation sparsity can improve language-model inference efficiency when specialized kernels skip computations associated with zero activations. Recent work has shown that sparsity can be exposed in pretrained models through post-hoc thresholding or induced during training, in both cases with measured inference speedups~\citep{liu2024teal,cetin2026sparser}. These results raise a practical design question: how should a researcher place and combine sparsification interventions to obtain more useful sparsity at an acceptable quality cost?

Existing approaches induce activation sparsity in different ways. TEAL applies post-hoc thresholding to pretrained models, ProSparse uses L1 regularization on FFN activations during training, and Q-Sparse applies top-k sparsification to both FFN and attention activations~\citep{liu2024teal,song2024prosparse,wang2024qsparse}. While each approach is effective in its own setting, it remains unclear how they compare under a common setup and how their effects interact when combined.% An intervention that works well in isolation may add little on top of another or incur a different quality cost when applied at additional sites, making sparsity alone insufficient to identify which choices drive the improvement.

In this work, we study three sparsification interventions: (i) \emph{pressure} through $\ell_1$ regularization on activations; (ii) \emph{thresholding}, i.e. mapping small activations exactly to zero; and (iii) \emph{placement}, which determines where these interventions are applied. To measure its effects, rather than report sparsity locally at individual activation sites, we define model-wide activation sparsity, $\Smodel$, as the fraction of forward matrix-multiplication products with a zero activation operand. We additionally derive an architecture-dependent sparsity ceiling, $\Smax$, corresponding to the maximum sparsity that can be achieved through the intervention set. Under matched pretraining conditions, we try to isolate the individual and interaction effects of each intervention and connect them to the amount of model-wide computation potentially avoided.

%% Comments:
%- Briefly justify the pythia-family choice (controlled recipe accross model sizes)
%- Do a list of findings, e.g.
%(i) thresholding seems to be more effective in generating sparsity, because the pressure creates smaller activation density mass, but without thresholing it does not convergen to zero altought we're using L1;
%(ii) FFNs seems more prone to sparsification. When applied the same thresholding and pressure, FFNs achieve a higher local sparsity then self-attention sites;
%(iii) accross model sizes, the more aggressive thresholding+pressure recipes achieve higher model-wide sparsity than post-hoc clipping over the baseline. But just changing the GeLU to ReLU (A1-H) improves de post-hoc clipping frontier.. which gets better and better with increased sizes.
% Paragraph evidence: introduction-findings.md; Analysis 020 observations/O003-cross-size.md.
We pretrain Pythia-family models from scratch for controlled comparisons across model sizes~\citep{biderman2023pythia}. Across 30 conditions at 14M, combining thresholding with pressure raises model-wide sparsity from $15.39\%$ to $27.48\%$, near its $29.95\%$ ceiling, for $+0.13$ validation loss. FFNs are sparser at moderate thresholds, while stronger thresholding also sparsifies attention. Selected combinations outperform aggressive post-hoc clipping at 70M and 410M. ReLU pretraining also improves subsequent clipping, with larger gains at larger sizes under the same clipping setting.

Finally, we use GPT-6-Astra~\citep{openaigpt6astra} to develop specialized kernels and test whether $\Smodel$ translates into faster inference. Across 42 iterations, the selected implementation reaches a $1.78\times$ full-model speedup over native PyTorch/SDPA on an RTX5090 for Pythia-14M at $27.48\%$ model-wide sparsity, near its $29.95\%$ ceiling. Across all 30 14M models, speedup is linearly correlated with $\Smodel$ ($R^2=0.817$). Mapping architectural and train-time interventions to speedup therefore requires kernel development, which coding agents can support~\citep{luo2026agentic}.
